\documentclass[12pt,a4paper,oneside]{article}
\usepackage[brazil]{babel}
\usepackage[utf8]{inputenc}
\usepackage{olymp}
\usepackage{color}
\usepackage[dvipsnames]{xcolor}
\usepackage{listings}

\setcounter{problem}{3}

\begin{document}

\begin{problem}{Número espelhado}{espelho}{2}

Nesta questão você deve escrever um programa que lê um número e o
escreve espelhado. Por exemplo, para a entrada 12345 deve ser
escrito 54321.

Existem vários maneiras de se fazer isto. Nesta questão, entretanto, o
número espelhado deve ser \underline{escrito por uma função recursiva}.

%\Notes
%
%Nesta questão, seu programa deve usar a função \texttt{main} dada acima exatamente como está, não a modifique! Sua tarefa é modificar a função \texttt{fatorial}.


\InputFile

A entrada contém um único número natural $N$. Restrições: $0 < N < 2000000000$.

\OutputFile

Seu programa deve gerar uma linha de saída contendo o número $N$ espelhado.

\Example


\begin{example}
\exmp{
12345
}{
54321
}%
\end{example}

\begin{example}
\exmp{
102030405
}{
504030201
}%
\end{example}


\begin{example}
\exmp{
1000000
}{
0000001
}%
\end{example}

\vspace{0.5cm}
DICA: use \% e /
%Comentários sobre o exemplo, se necessário.

\end{problem}

\end{document}
